\scanpage{59}
\mspage{52}

aura, pour $I$ dans $\mathrm{Cat}$

\[ \begin{aligned}
M(I)&=K^{\unsure{\bullet|\cdot}}(A(I))\\
W(I)&=\underbrace{W_{A(I)}}_{\text{quasi-isomorphismes}}\unsure{\subset|\in}\mathrm{Fl}(K^\bullet(A(I)))
\end{aligned} \]

et de même pour $M^\alpha(I),W^\alpha(I)$ $(\alpha\in\{+,-,b\})$, du moins \uline{si $I$ est fini}. Il y aussi les sous-catégories $\unsure{K^{\geq0}|K^{>0}}(A)$ et $K^{\leq0}(A)$ des complexes de \unsure{cochaines|cochaînes} resp. de \unsure{chaines|chaînes} au sens strict (avec $K^i=0$ si $i\geq0$, resp\unsure{t}. si $i\leq0$). On aura pour $I$ quelconque

\[ \left\{\begin{aligned}
K^{\geq0}(A)(I)&\simeq K^{\geq0}(A(I))\\
W_A^{\geq0}(I)&=W_{A(I)}^{\geq0}
\end{aligned}\right. \]

et de $=$ pour $K^{\leq0},W^{\leq0}$.

D'après Quillen \unsure{pag.|p.} 1.2, $(K^+(A),W_A^+)$ est quillénisable si $A$ a assez d'injectifs, $(K^-(A),W_A^-)$ l'est si $A$ a assez de projectifs. Il \unsure{ne dit|n'en est} rien pour $(K^\bullet(A),W_A^{\unsure{*|\bullet}})$ sans plus -- pourtant on fait de l'algèbre cohomologique avec \deleted{$K(A),W_A)$} $D^{\unsure{?|\bullet}}(A)$ sans plus, et $=$ sans faire d'hypothèse d'existence d'injectifs ou de projectifs. J'ai un peu oublié ce que Verdier a fait à ce sujet -- je devrais quand même regarder sa thèse (!), ou bien l'exposé de Hartshorne,
% archive sheet number: 58
